% GENERATED FILE. Edit canonical lessons, then run npm run generate:books.
% canonical-source-hash: fnv1a64:4fd55dbedbd44993
% canonical-lessons: FR-C111-tapis, FR-C111-television, FR-C111-ordinateur, FR-C111-refrigerateur, FR-C111-radio

\chapter{Rug, Television, Computer, Fridge, Radio}
\label{ch:fr-rug-television-computer-fridge-radio}
\hlchaptermodality{111}

\begin{chapteropening}
\textbf{By the end of this chapter:} \emph{I can name a rug, a television, a computer, a fridge, a radio.}

Complete the last lesson of chapter 111: I can name a rug, a television, a computer, a fridge, a radio.
\end{chapteropening}

\section[le tapis]{le tapis — a rug, a carpet}

\label{lesson:FR-C111-tapis}



\begin{quote}
Before the new one: say the French for a journey, then the French for holidays.
\end{quote}

\subsection*{You'll want to know: le tapis}
\textbf{le tapis} — \textquotedblleft{}a rug, a carpet\textquotedblright{}.

\begin{cousinweb}
\textbf{tapis} $\leftarrow$ Old French \emph{tapiz} $\leftarrow$ Byzantine Greek \emph{tapetion}, from Greek \emph{tapes}, \textquotedblleft{}a carpet\textquotedblright{} — borrowed. (EtymDB 2.1)
\end{cousinweb}

\begin{grammarlens}[title={What you've built}]
One more word for this chapter.
\end{grammarlens}

\subsection*{Your turn}
\begin{itemize}
\raggedright
  \item \emph{Say it:} \emph{le tapis}
  \item \emph{Say it:} \emph{le tapis}, once more
  \item \emph{Say it:} say \emph{le tapis}
  \item \emph{Recall:} say the French for a journey, then the French for holidays, then say \emph{le tapis} again
\end{itemize}

\begin{tcolorbox}[breakable,colback=teal!4,colframe=teal!35!black,title={Before you move on}]
What does le tapis mean? (\textquotedblleft{}A rug, a carpet\textquotedblright{}.) Say it once more.
\end{tcolorbox}
\section[la télévision]{la télévision — a television}

\label{lesson:FR-C111-television}



\begin{quote}
Before the new one: say the French for holidays, then the French for a rug.
\end{quote}

\subsection*{You'll want to know: la télévision}
\textbf{la télévision} — \textquotedblleft{}a television\textquotedblright{}.

\begin{cousinweb}
\textbf{télévision} — a 20th-century coinage from Greek \emph{tele}, \textquotedblleft{}far off\textquotedblright{}, and \emph{vision}, $\leftarrow$ Latin \emph{visio}, \textquotedblleft{}seeing\textquotedblright{} (Lewis \& Short).
\end{cousinweb}

\begin{grammarlens}[title={What you've built}]
One more word for this chapter.
\end{grammarlens}

\subsection*{Your turn}
\begin{itemize}
\raggedright
  \item \emph{Say it:} \emph{la télévision}
  \item \emph{Say it:} \emph{la télévision}, once more
  \item \emph{Say it:} say \emph{la télévision}
  \item \emph{Recall:} say the French for holidays, then the French for a rug, then say \emph{la télévision} again
\end{itemize}

\begin{tcolorbox}[breakable,colback=teal!4,colframe=teal!35!black,title={Before you move on}]
What does la télévision mean? (\textquotedblleft{}A television\textquotedblright{}.) Say it once more.
\end{tcolorbox}
\section[l'ordinateur]{l'ordinateur — a computer}

\label{lesson:FR-C111-ordinateur}



\begin{quote}
Before the new one: say the French for a rug, then the French for a television.
\end{quote}

\subsection*{You'll want to know: l'ordinateur}
\textbf{l'ordinateur} — \textquotedblleft{}a computer\textquotedblright{}.

\begin{cousinweb}
\textbf{ordinateur} $\leftarrow$ Latin \emph{ordinator}, \textquotedblleft{}one who sets in order\textquotedblright{} (EtymDB 2.1) — a learned borrowing, chosen in 1955 as the French name for the computer.
\end{cousinweb}

\begin{grammarlens}[title={What you've built}]
One more word for this chapter.
\end{grammarlens}

\subsection*{Your turn}
\begin{itemize}
\raggedright
  \item \emph{Say it:} \emph{l'ordinateur}
  \item \emph{Say it:} \emph{l'ordinateur}, once more
  \item \emph{Say it:} say \emph{l'ordinateur}
  \item \emph{Recall:} say the French for a rug, then the French for a television, then say \emph{l'ordinateur} again
\end{itemize}

\begin{tcolorbox}[breakable,colback=teal!4,colframe=teal!35!black,title={Before you move on}]
What does l'ordinateur mean? (\textquotedblleft{}A computer\textquotedblright{}.) Say it once more.
\end{tcolorbox}
\section[le réfrigérateur]{le réfrigérateur — a fridge}

\label{lesson:FR-C111-refrigerateur}



\begin{quote}
Before the new one: say the French for a television, then the French for a computer.
\end{quote}

\subsection*{You'll want to know: le réfrigérateur}
\textbf{le réfrigérateur} — \textquotedblleft{}a fridge\textquotedblright{}.

\begin{cousinweb}
\textbf{réfrigérateur} — formed from the verb for \textquotedblleft{}to chill\textquotedblright{}, a learned borrowing of Latin \emph{refrigerare}, \textquotedblleft{}to make cool\textquotedblright{}, from \emph{frigus}, \textquotedblleft{}cold\textquotedblright{} (Lewis \& Short).
\end{cousinweb}

\begin{grammarlens}[title={What you've built}]
One more word for this chapter.
\end{grammarlens}

\subsection*{Your turn}
\begin{itemize}
\raggedright
  \item \emph{Say it:} \emph{le réfrigérateur}
  \item \emph{Say it:} \emph{le réfrigérateur}, once more
  \item \emph{Say it:} say \emph{le réfrigérateur}
  \item \emph{Recall:} say the French for a television, then the French for a computer, then say \emph{le réfrigérateur} again
\end{itemize}

\begin{tcolorbox}[breakable,colback=teal!4,colframe=teal!35!black,title={Before you move on}]
What does le réfrigérateur mean? (\textquotedblleft{}A fridge\textquotedblright{}.) Say it once more.
\end{tcolorbox}
\section[la radio]{la radio — a radio}

\label{lesson:FR-C111-radio}



\begin{quote}
Before the new one: say the French for a computer, then the French for a fridge.
\end{quote}

\subsection*{You'll want to know: la radio}
\textbf{la radio} — \textquotedblleft{}a radio\textquotedblright{}.

\begin{cousinweb}
\textbf{radio} — short for \emph{radiotélégraphie}, built on Latin \emph{radius}, \textquotedblleft{}a rod, a ray\textquotedblright{} (Lewis \& Short).
\end{cousinweb}

\begin{grammarlens}[title={What you've built}]
One more word for this chapter.
\end{grammarlens}

\subsection*{Your turn}
\begin{itemize}
\raggedright
  \item \emph{Say it:} \emph{la radio}
  \item \emph{Say it:} \emph{la radio}, once more
  \item \emph{Say it:} say \emph{la radio}
  \item \emph{Recall:} say the French for a computer, then the French for a fridge, then say \emph{la radio} again
\end{itemize}

\begin{tcolorbox}[breakable,colback=teal!4,colframe=teal!35!black,title={Before you move on}]
What does la radio mean? (\textquotedblleft{}A radio\textquotedblright{}.) Say it once more.
\end{tcolorbox}
